\documentclass[11pt]{article}

% Packages
\usepackage[margin=1in]{geometry}
\usepackage{amsmath, amssymb}
\usepackage{enumitem}
\usepackage{titlesec}
\usepackage{hyperref}
\usepackage{xcolor}

% Formatting
\setlength{\parskip}{6pt}
\setlength{\parindent}{0pt}

\titleformat{\section}{\large\bfseries}{\thesection.}{0.5em}{}
\titleformat{\subsection}{\normalsize\bfseries}{\thesubsection.}{0.5em}{}

% Custom commands
\newcommand{\prdtitle}[1]{\begin{center}\LARGE\bfseries #1\end{center}}
\newcommand{\prdauthor}[1]{\begin{center}\normalsize #1\end{center}}
\newcommand{\prdmeta}[1]{\begin{center}\small\textcolor{gray}{#1}\end{center}}

\begin{document}

% Title Block
\prdtitle{Rubric-Guided LLM Grading with Continuous Monitoring}
\prdauthor{Author: Miguel Bueno}
\prdmeta{Last Updated: \today}

\vspace{1em}

% -------------------------
\section{Introduction}
\textbf{Overview:} This product is a rubric-guided grading system that uses a
large language model (LLM) to evaluate free-response assignments at scale.
Instructors define structured rubrics that specify grading criteria, and the
system applies those rubrics to student responses to generate scores and
feedback.

The product operates in the context of higher education and other instructional
settings where open-ended responses are common but costly to evaluate. It is
designed to support instructors by reducing grading workload while maintaining
transparency and alignment with clearly defined evaluation standards.

\textbf{Background:} 
Free-response assignments are widely used in instructional settings because they
assess higher-order skills such as reasoning, synthesis, and written communication.
However, they are significantly more time-consuming to evaluate than structured
formats. As a result, instructors often face a tradeoff between assigning richer,
open-ended work and managing grading workload.

This constraint directly impacts students. Because grading is time-intensive,
feedback is often delayed by several days, reducing its usefulness for learning and
limiting opportunities for iteration and improvement.

In practice, grading is also subject to \textbf{intra-rater variability}, where the
same instructor may apply criteria inconsistently over time. In addition, grading
can be influenced by \textbf{implicit biases}, such as prior expectations about a
student’s performance or writing style, which may affect scores beyond the content
of the response. These factors can reduce the consistency and perceived fairness of
evaluation.

This creates a need for a system that can scale grading of open-ended responses
while maintaining consistency, transparency, and alignment with explicitly defined
evaluation criteria, while reducing bias and improving feedback timeliness for students.

\textbf{Scope:}

\textbf{In-scope}
\begin{enumerate}[leftmargin=*]
    \item \textbf{Rubric creation and management} \\
    Structured definition of grading criteria, including dimensions, descriptions, and scoring levels
    
    \item \textbf{LLM-based grading of free-response text} \\
    Automated evaluation of student responses with dimension-level scores and rubric-aligned justifications
    
    \item \textbf{Batch processing of student submissions} \\
    Scalable grading workflows supporting assignment- and class-level processing
    
    \item \textbf{Instructor monitoring via sampling} \\
    Review of randomly and selectively sampled graded responses to assess and maintain quality
    
    \item \textbf{Quality tracking and basic diagnostics} \\
    Estimation of grading error rates and identification of low-confidence or inconsistent outputs
    
    \item \textbf{Student-facing feedback} \\
    Delivery of rubric-aligned scores and explanations to support student understanding
    
    \item \textbf{Privacy and data protection controls} \\
    Handling of student data in compliance with applicable regulations (e.g., FERPA), including minimization, restricted access, and anonymization where appropriate
\end{enumerate}

\textbf{Out-of-scope}
\begin{enumerate}
    \item \textbf{Full automation without instructor oversight} \\
    Fully autonomous grading workflows without any instructor review or monitoring

    \item \textbf{Learning individual instructor grading behavior} \\
    Calibration to instructor-specific labels or imitation of individual grading patterns

    \item \textbf{Multi-instructor consensus or cross-course standardization} \\
    Harmonization of grading across instructors, courses, or institutions

    \item \textbf{Advanced LMS integrations (Canvas, etc.)} \\
    Deep platform integrations beyond basic export and import functionality

    \item \textbf{Non-text or multimodal grading} \\
    Evaluation of diagrams, code execution, or other non-textual formats in the initial version

    \item \textbf{Use of student data for unrelated model training} \\
    Use of student submissions for training general-purpose models without explicit consent
\end{enumerate}

% -------------------------
\section{Objectives}
\textbf{Goals:}
\begin{enumerate}[leftmargin=*]
    \item \textbf{Primary objective}
    \begin{enumerate}
        \item \textbf{Scalable, rubric-aligned grading of free-response work} \\
        Consistent application of instructor-defined evaluation criteria across large volumes of student responses
    \end{enumerate}

    \item \textbf{Secondary objectives}
    \begin{enumerate}
        \item \textbf{Reduced instructor grading burden} \\
        Decreased time spent on manual grading through automated evaluation and targeted review

        \item \textbf{Improved feedback timeliness for students} \\
        Faster return of meaningful, rubric-based feedback to support learning and iteration

        \item \textbf{Increased consistency and transparency in grading} \\
        More uniform application of grading criteria with clear justification tied to rubric dimensions

        \item \textbf{Detection and mitigation of potential bias} \\
        Identification of systematic deviations from rubric-based evaluation that may reflect implicit bias

        \item \textbf{Establishment of a measurable quality control loop} \\
        Ongoing estimation of grading accuracy through sampling-based monitoring

        \item \textbf{Privacy, security, and regulatory compliance} \\
        Protection of student PII and adherence to applicable legal and institutional requirements
    \end{enumerate}
\end{enumerate}

% -------------------------

\textbf{Success Metrics:}
\begin{itemize}[leftmargin=*]
    \item Metric 1:
    \item Metric 2:
    \item Guardrails:
\end{itemize}

% -------------------------
\section{Stakeholders}
\begin{itemize}[leftmargin=*]
    \item \textbf{Product Owner:}
    \item \textbf{Engineering:}
    \item \textbf{Data Science:}
    \item \textbf{Design:}
    \item \textbf{Other:}
\end{itemize}

% -------------------------
\section{Use Cases}
\begin{enumerate}[leftmargin=*]
    \item \textbf{User Story:} As a \underline{\hspace{2cm}}, I want to \underline{\hspace{4cm}} so that \underline{\hspace{4cm}}.
    
    \textbf{Acceptance Criteria:}
    \begin{itemize}
        \item 
        \item 
    \end{itemize}

    \item \textbf{User Story:} As a \underline{\hspace{2cm}}, I want to \underline{\hspace{4cm}} so that \underline{\hspace{4cm}}.
    
    \textbf{Acceptance Criteria:}
    \begin{itemize}
        \item 
        \item 
    \end{itemize}
\end{enumerate}

% -------------------------
\section{Aspects}
\subsection{Functional Requirements}
\begin{itemize}[leftmargin=*]
    \item 
    \item 
\end{itemize}

\subsection{Non-Functional Requirements}
\begin{itemize}[leftmargin=*]
    \item Performance:
    \item Reliability:
    \item Scalability:
    \item Privacy / Safety:
\end{itemize}

\subsection{Data Requirements}
\begin{itemize}[leftmargin=*]
    \item Data sources:
    \item Logging:
    \item Metrics instrumentation:
\end{itemize}

% -------------------------
\section{Open Questions}
\begin{itemize}[leftmargin=*]
    \item 
    \item 
\end{itemize}

% -------------------------
\section{Milestones}
\begin{itemize}[leftmargin=*]
    \item \textbf{Phase 1:} Description --- Date
    \item \textbf{Phase 2:} Description --- Date
    \item \textbf{Launch:} Date
\end{itemize}

\end{document}
